\newpage{}
\section{Data}
\label{sec:data}
\subsection[Input Data]{Input Data $(\tensor{x}, {}_{x}\tensor{m})$}
\label{subsec:input-data}

The raw market observations are a large and structured window of
normalized \texttt{kline} sequences: several channels, several points, and several
features per edge. As introduced in Section~\ref{sec:edges-and-nodes}, 
$\tensor{x}_{u\,v}(t)$ is the normalized%
\textsuperscript{\textcolor{FootnoteField}{\ref{fn:data-normalized}}}
edge-payload tensor attached to the directed edge
$e=(u,v)\in\mathcal{E}$ at time $t$.
\[
  \tensor{x}_{u\,v}(t)
  \in \R^{C \times H_x \times D_x},
\]
Where $C$ indexes the channels%
\note{
  Channels will use different resolutions or time scales,
  such as \field{1m}, \field{1h}, or \field{1d},... for minute, hour, or day
  intervals. The whole list of time scales is:
  \field{1w}, \field{3d}, \field{1d}, \field{12h}, \field{8h},
  \field{6h}, \field{4h}, \field{2h}, \field{1h}, \field{30m},
  \field{15m}, \field{5m}, \field{3m} and \field{1m}. We don't try to model high frequency. 
}, 
$H_x$ is the input sequence length%
\note{
  Each channel may contribute a different native input length, 
  so the dataloader pads them to a common length to preserve a rectangular tensor shape 
  and records the pad into the mask so the effective sequence length is not altered.
  The shared time anchor is the last available record on each channel not after the
  common anchor time.  Positions across channels therefore have the same role
  relative to the anchor, but they are not required to have identical raw
  \texttt{close\_time} feature on the last element of the sequence.  See
  Appendix~\ref{app:sequence-synchronization} for more details.
}, 
and $D_x$ is the model-feature width%
\note{
  The stored \field{kline\_t} fields are, in order:
  \field{open\_time}, \field{open\_price}, \field{high\_price},
  \field{low\_price}, \field{close\_price}, \field{volume},
  \field{close\_time}, \field{quote\_asset\_volume},
  \field{number\_of\_trades}, \field{taker\_buy\_base\_volume}, and
  \field{taker\_buy\_quote\_volume}.  The model feature vector uses the four
  price fields and five activity fields; \field{open\_time} and
  \field{close\_time} are alignment keys rather than normalized model
  features.
}.
\refstepcounter{footnote}\label{fn:data-normalized}\footnotetext{Here
\textit{normalized} means that raw candle fields have been mapped to
model coordinates.  Each OHLC price $p(t)$ is represented by the log-return
$\log\!\bigl(p(t)/\mathrm{close\_price}(t-1)\bigr)$, where
\field{close\_price}$(t-1)$ is the previous close on the same channel.
Non-price activity fields are mapped coordinatewise by
$\operatorname{log1p}(u)=\log(1+u)$.  For more details see
Appendix~\ref{app:normalization}.} 
The sample validity mask%
\note{
  Market data are not a complete mathematical grid.  Source records may be
  missing, duplicated, stale, out of order, or unusable after alignment; the
  normalization step can also be undefined when a required reference value, such
  as the previous valid close, is unavailable.  The mask marks these positions as
  invalid instead of forcing fabricated feature values into the tensor.
} tensor is:
\[
  {}_{x}\tensor{m}_{u\,v}(t) \in \{0,1\}^{C \times H_x},
\]
where ${}_{x}m_{u\,v}^{c\,i}(t)=1$ means that feature slice
$\tensor{x}_{u\,v}^{c\,i}(t)$ is valid.\\

For one sample, one channel, and one input position, the normalized candle slice is%
\note{
  The tuple $(c,i)$ selects the active channel $c$, and input sequence position
  $i$, while the lower label $u\,v$ keeps the edge carried by the sample
  visible, with $e=(u,v)$.  $\tensor{x}_{u\,v}^{c\,i}$ denotes the full
  $D_x$-dimensional normalized candle vector; an individual scalar entry would
  be indexed as $x_{u\,v}^{c\,i\,d}$.
  Note that $i$ indexes the channel's sequence, 
  not to be confused with $t$ that index anchor that synchronizes all channel's sequences.
}
\[
  \tensor{x}_{u\,v}^{c\,i}(t) \in \R^{D_x},
  \qquad
  {}_{x}\tensor{m}_{u\,v}^{c\,i}(t) \in \{0,1\},
\]

For batched notation%
\note{
  Batching adds the leading sample axis, giving
  $(\tensor{x}_{u\,v}^{b}(t))_{b=1}^{B}
  \in\R^{B \times C \times H_x \times D_x}$. Same for the mask 
  $({}_{x}\tensor{m}_{u\,v}^{b}(t))_{b=1}^{B}
  \in\R^{B \times C \times H_x}$
}, stack the window with an extra rank as
\[
  \tensor{x}(t)
  \in \R^{B \times C \times H_x \times D_x},
  \qquad
  {}_{x}\tensor{m}(t)
  \in \{0,1\}^{B \times C \times H_x}.
\]











\newpage{}

\subsection[Targets for supervised learning]{Targets for supervised learning ($\tensor{f}$, ${}_{f}\tensor{m}$)}%
\label{subsec:targets-for-supervised-learning}


Given an anchor time $t$, the supervised target is a normalized%
\note{
  Future windows use the same log-return and \(\operatorname{log1p}\) normalization
  rules as the input tensor; see footnote~\ref{fn:data-normalized} and
  Appendix~\ref{app:normalization}.
}
future window%
\note{
  In the dataloader, past and future windows are organized around the same
anchor time $t$.  The input window $\tensor{x}_{u\,v}(t)$ is back-aligned: for each
channel, its last valid record is the anchor record at or immediately before
the common anchor key, and any channel padding is added at the front.  The
future window $\tensor{f}_{u\,v}(t)$ is forward-aligned: for each channel, its first
valid record is the first record after that anchor, and any channel padding is
added at the end.  Thus the two tensors are aligned on opposite sides of the
same anchor rather than by a row-by-row timestamp identity across channels.
See Appendix~\ref{app:sequence-synchronization}.
}. 
\note{
  The future window is the part of the same edge-indexed kline signal that
lies after the input window ending at time $t$.  It is supervised because it is
known during training and hidden during inference.
}.
\[
  \tensor{f}_{u\,v}(t) \in \R^{C \times H_f \times D_f},
\]
with validity mask%
\note{
  The future mask plays the same role as the input mask: it excludes target
positions that are unavailable or not meaningful after alignment and
normalization, so loss terms are computed only on valid future targets.
}:
\[
  {}_{f}\tensor{m}_{u\,v}(t) \in \{0,1\}^{C \times H_f}.
\]
For one channel and one forecast position, the normalized future target slice is:
\[
  \tensor{f}_{u\,v}^{c\,i}(t) \in \R^{D_f}.
\]
Here $H_f$ is the forecast sequence length and $D_f$ is the number of target
coordinates. \\

For batched notation, stack targets as%
\note{
  For batching operations, simply stack the feature and mask tensors as: 
  $\left(\tensor{f}_{u\,v}^{b}(t)\right)_{b=1}^{B} \in \R^{B \times C \times H_f \times D_f}$, 
  and 
  $\left({}_{f}\tensor{m}_{u\,v}^{b}(t)\right)_{b=1}^{B} \in \{0,1\}^{B \times C \times H_f}$.
}
\[
  \tensor{f}(t) \in \R^{B \times C \times H_f \times D_f},
  \qquad
  {}_{f}\tensor{m}(t) \in \{0,1\}^{B \times C \times H_f}.
\]

\calloutbox{Non-Forecasting Claims \& Future Uncertainty Disambiguation%
\note{
  This document is a technical description of a research system.  It is
provided for scientific and engineering discussion only and does not
constitute investment advice, a trading recommendation, or an offer to buy or
sell any financial instrument.  Model outputs are uncertain, may be wrong, and
do not guarantee future results; any operational, financial, or regulatory use
remains the sole responsibility of the user.
}.}{\\%
We do not claim deterministic knowledge of future market values.  Rather, the
future window is used as supervised data for estimating conditional
probability, from which quantities such as expected return and risk may be
read.  The model represents uncertainty about future outcomes; it does not
remove that uncertainty.  See Section~\ref{sec:edge-forecast-head} for the
density-mixture formulation.
}

See section~\ref{sec:representation} for how we use the same embedding network to 
encode the data on any edge to a unified representation.











\newpage{}
\subsection[Portfolio Holding State]{Portfolio Holding States $(\tensor{h}, \tensor{\eta})$}
\label{subsec:exposure}

Market observations ($\tensor{x}$) and portfolio holdings ($\tensor{h}$) live on different objects of the
graph.  Observations are on the edges, holdings%
\note{
  Holdings $\tensor{h}(t)$ represent the investment positions at any anchor-time $t$.
} 
are on the nodes. Nodes also have valuation $\tensor{\eta}$%
\note{
  We use machine learning to derive $\tensor{\eta}(t)$ 
  and estimate the valuation of one asset relative to another.
}.

At anchor time $t$, the amount of units we hold for asset $u$ is%
\note{
  When the nodes are concrete assets, ${h}_u(t)$ may be read as signed inventory
  in the native unit of asset $u$.  More generally, it is the quantity recorded in the portfolio records, 
  we do this for all nodes (assets) even if we have zero of such asset: 
  $\tensor{h}(t) = \bigl({h}_u(t)\bigr)_{u\in\mathcal{V}} \in \R^N$
}: 
\[
  u\in\mathcal{V}
  \quad
  \longmapsto
  \quad
  {h}_u(t)\in\R,
\]
In general, to reference all assets at once we say: 
\[
  \tensor{h}(t)\in\R^{N}.
\]

The sign%
\note{
  A positive value denotes a long holding, claim, or positive exposure to node
  $u$.  A negative value denotes a short position, borrowed inventory,
  liability, or any construction whose economic direction is opposite to holding
  the node.  The value ${h}_u(t)=0$ means that the portfolio has no active holding
  at node $u$ at time $t$.
} 
of $h_u(t)$ records whether the holding is long or short%
\note{
  No normalization is imposed for this variable. In particular, 
  $\tensor{h}(t)$ is not a probability vector, not a vector of
  portfolio weights, and not required to have nonnegative entries 
  or to sum to one.
  Its coordinates also need not share a common external denomination
  before an accounting map is chosen.
}.

At the nodes we attach a separate object, the node valuation $\tensor{\eta}_u(t)$.%
\note{
  $\tensor{\eta}$ represents a node-level valuation coordinate, recovered through edge relations.
}
It is used to assign relative value and compare assets with and against each other.
\[
  \tensor{\eta}(t)\in\R^{N}.
\]
